\subsection{Sprint 2}

\subsubsection*{Desenvolvimento da US-9 --- Visualizar e Editar Perfil do Aluno}

O desenvolvimento da funcionalidade \textbf{US-9 --- Visualizar e Editar Perfil do Aluno}
foi iniciado logo após sua disponibilização para implementação. A interface já possuía
uma proposta previamente elaborada no Figma, contendo campos semelhantes aos
utilizados no cadastro de usuários do sistema.

Durante a análise dos requisitos, surgiram dúvidas relacionadas aos dados que deveriam
estar disponíveis para edição no perfil do aluno. A documentação da tarefa não
apresentava detalhes suficientes para definir com clareza quais informações seriam
prioritárias para alteração pelos usuários. Para solucionar essa questão, foi realizado
contato com a equipe \ac{ages}, que forneceu orientações complementares e uma lista
dos campos necessários para a implementação da funcionalidade.

Essa situação evidenciou uma dificuldade que vinha sendo observada em algumas tasks
do projeto. Em determinados casos, as descrições registradas no ClickUp apresentavam
nível reduzido de detalhamento, abrindo espaço para diferentes interpretações dos
requisitos e ocasionando retrabalho durante o desenvolvimento.

Diante desse cenário, a estratégia adotada consistiu em iniciar as atividades o mais
cedo possível após sua publicação. Essa abordagem permitiu identificar dúvidas e
inconsistências ainda nas etapas iniciais do desenvolvimento, contribuindo para a
redução de retrabalho e para uma maior assertividade na implementação das
funcionalidades.

\subsubsection*{Resultados da Sprint}

Durante esta sprint, a Squad Ninja recebeu avaliações positivas em relação à qualidade
das entregas realizadas. As funcionalidades desenvolvidas apresentaram baixo número de
erros e pendências durante o processo de revisão, facilitando a conferência e a
integração por parte dos demais integrantes da equipe.

Como consequência, o processo de incorporação das funcionalidades ao sistema ocorreu
de forma mais eficiente, exigindo poucas correções após a conclusão do
desenvolvimento. Esse resultado demonstrou um bom alinhamento entre planejamento,
implementação e validação das entregas.

Na retrospectiva da sprint, foram registrados cinco reconhecimentos (\textit{kudos})
relacionados à participação ativa nas atividades do projeto e ao apoio prestado durante
o processo de integração das funcionalidades. Além disso, a equipe destacou a qualidade
da apresentação realizada pelos \ac{ages} IV, tanto pela organização do conteúdo quanto
pela qualidade visual empregada na exposição dos resultados.

Outro aspecto positivo apontado foi a evolução da dinâmica de trabalho do quarteto,
que apresentou maior colaboração, comunicação e organização ao longo da sprint. As
entregas realizadas permaneceram alinhadas ao planejamento estabelecido, e as poucas
pendências existentes já haviam sido previamente identificadas e consideradas pela
equipe durante o acompanhamento das atividades.

\subsubsection*{Ponto de Atenção da Sprint}

O principal desafio identificado durante a sprint esteve relacionado ao alinhamento
entre as atividades de frontend e backend. Em alguns casos, a divisão dessas tarefas
entre diferentes responsáveis resultou em interpretações distintas dos requisitos,
ocasionando dificuldades durante a integração das funcionalidades.

Embora os componentes desenvolvidos estivessem tecnicamente corretos, algumas
implementações não atendiam integralmente às necessidades previstas para a
funcionalidade, devido à falta de detalhamento dos requisitos de integração entre as
partes do sistema.

Como ação de melhoria, foi acordado que as definições das tasks passarão a apresentar
maior nível de detalhamento, especialmente em relação às regras de negócio e aos
requisitos de comunicação entre frontend e backend. Além disso, sempre que possível,
haverá um acompanhamento mais integrado entre os responsáveis pela definição dessas
atividades, reduzindo ambiguidades e facilitando o desenvolvimento das funcionalidades.

A expectativa é que essas medidas contribuam para diminuir problemas de integração,
aumentar a clareza dos requisitos e tornar o processo de desenvolvimento mais eficiente
nas próximas sprints.
